\chapter{Entity--Relationship (ER) Modeling}

\section{Identification of Entities}
The model has 30 entity types in six modules. Each becomes one table (Chapter 4); the two M:N and
ternary relationships with attributes become the other two tables, and \tbl{schema\_migrations} records applied
schema versions and is not part of the domain. Table~\ref{tab:entities} lists every entity with its identifier;
the diagrams show each entity's principal attributes, and Section~4.3 lists all 342 columns.
Four entities are \textbf{weak} (Section~3.4).

{\scriptsize
\begin{longtable}{@{}llllp{4.2cm}@{}}
\caption{Entities, their module, table and identifier}\label{tab:entities}\\
\toprule
\textbf{Entity} & \textbf{Mod.} & \textbf{Table} & \textbf{Type} & \textbf{Identifier (key / partial key)} \\\midrule\endfirsthead
\toprule\textbf{Entity} & \textbf{Mod.} & \textbf{Table} & \textbf{Type} & \textbf{Identifier} \\\midrule\endhead
\rows{figures/er/entity-table.tex}
\bottomrule
\end{longtable}}

\section{Identification of Relationships}
There are 69 relationship types. 64 are binary 1:N, three are binary 1:1 (\textsc{becomes},
\textsc{yields}, \textsc{is\_doctor}) and two carry attributes: \textsc{filled\_from} is M:N (a prescription line can be filled
from several batches, and a batch fills many lines) and \textsc{occupies} is ternary (bed, admission, patient).
Together they account for every one of the 72 relationship foreign keys of the schema; a script checks this,
and checks each participation and cardinality below against the column's nullability and unique keys.
In Table~\ref{tab:rels} the ratio and participation are given in the order the entities are listed
(T = total, P = partial); $^{\ast}$ marks an identifying relationship.

{\tiny\setlength{\tabcolsep}{2.2pt}
\begin{longtable}{@{}llll|llll@{}}
\caption{Relationships, cardinality ratios and participation}\label{tab:rels}\\
\toprule
\textbf{Relationship} & \textbf{Entities} & \textbf{Ratio} & \textbf{Part.} & \textbf{Relationship} & \textbf{Entities} & \textbf{Ratio} & \textbf{Part.} \\\midrule\endfirsthead
\toprule\textbf{Relationship} & \textbf{Entities} & \textbf{Ratio} & \textbf{Part.} & \textbf{Relationship} & \textbf{Entities} & \textbf{Ratio} & \textbf{Part.} \\\midrule\endhead
\rows{figures/er/rel-table.tex}
\bottomrule
\end{longtable}}

\section{ER Diagram}
The diagrams use Chen notation (legend in Figure~\ref{fig:legend}). Figure~\ref{fig:overview} shows the six
modules and the 39 relationships between modules. Figures~\ref{fig:m1a}--\ref{fig:m6b} show the modules one
page at a time; three modules take two pages. An entity drawn in another module appears as a dashed reference box
(``see Module $n$''), and every relationship between two modules is drawn in both, with the same name and cardinality.

\begin{figure}[p]
\centering
% Full legend of the Chen notation used in Figures 3.2 to 3.11.
\begin{tikzpicture}[x=1cm,y=1cm, every node/.append style={}]
  \def\row#1#2#3{\node[anchor=west, font=\footnotesize] at (#1+2.35, #2) {#3};}
  % column 1
  \node[ent] at (1,0) {PATIENT};                         \row{0}{0}{Entity (strong): own primary key}
  \node[weakent] at (1,-0.9) {BED};                      \row{0}{-0.9}{Weak entity: identified only with its owner}
  \node[refent] at (1,-1.85) {DOCTOR\\[-1pt]{\tiny see Module 6}}; \row{0}{-1.85}{Entity drawn in another module (reference)}
  \node[rel] at (1,-2.85) {BOOKS};                       \row{0}{-2.85}{Relationship}
  \node[idrel] at (1,-3.8) {HAS\_BED};                   \row{0}{-3.8}{Identifying relationship of a weak entity}
  \node[att] at (1,-4.6) {name};                         \row{0}{-4.6}{Attribute}
  \node[att] at (1,-5.3) {\underline{patient\_id}};      \row{0}{-5.3}{Key attribute (underlined)}
  % column 2
  \node[att] at (9.0,0) {\dashuline{bed\_number}};       \row{8.4}{0}{Partial key of a weak entity (dashed underline)}
  \node[att] (c) at (9.0,-0.75) {blood\_pressure};
  \node[att] (c1) at (8.3,-1.45) {bp\_systolic}; \node[att] (c2) at (9.7,-1.45) {bp\_diastolic}; \draw (c)--(c1) (c)--(c2);
                                                          \row{8.4}{-1.05}{Composite attribute and its components}
  \node[att,double,double distance=0.8pt] at (9.0,-2.25) {allergies}; \row{8.4}{-2.25}{Multivalued attribute (double ellipse)}
  \node[att,dashed] at (9.0,-3.0) {age};                 \row{8.4}{-3.0}{Derived attribute (dashed ellipse)}
  \draw[total] (8.2,-3.75) -- (9.8,-3.75);              \row{8.4}{-3.75}{Total participation: every instance takes part}
  \draw[partial] (8.2,-4.5) -- (9.8,-4.5);              \row{8.4}{-4.5}{Partial participation: some instances take part}
  \draw[partial] (8.2,-5.3) -- node[card,pos=0.2]{1} node[card,pos=0.8]{N} (9.8,-5.3); \row{8.4}{-5.3}{Cardinality written at each end: 1:1, 1:N, M:N}
\end{tikzpicture}

\caption{Legend of the Chen notation used in all ER diagrams}\label{fig:legend}
\vspace{6mm}
% Overview: the six modules and every relationship between two modules (labels: figures/er/overview-labels.tex,
% generated from gen/er-model.js). Module 2 sits inside the triangle 1-5-6 and Module 4 outside it, so no
% two module links cross.
\newcommand{\ovlAB}{PRESCRIBED 1:N\\RX\_IN 1:N\\RX\_FROM 1:N\\SURGERY\_OF 1:N\\SURGERY\_IN 1:N}
\newcommand{\ovlAD}{TESTED 1:N\\LAB\_IN 1:N\\LAB\_FROM 1:N}
\newcommand{\ovlAE}{STAY\_VISIT 1:N\\TAKEN\_ON 1:N\\STAY\_NOTE 1:N\\ADMITTED 1:N\\OCCUPIES N:M:1\\ARRIVES 1:N\\BILLED 1:N}
\newcommand{\ovlAF}{ATTENDS 1:N\\CONDUCTS 1:N\\OPENS 1:N\\RECORDS 1:N\\WRITES 1:N}
\newcommand{\ovlBC}{OF\_DRUG 1:N\\FILLED\_FROM M:N}
\newcommand{\ovlBE}{STAY\_RX 1:N\\STAY\_DOSE 1:N}
\newcommand{\ovlBF}{PRESCRIBES 1:N\\GIVEN\_BY 1:N\\REQUESTS 1:N}
\newcommand{\ovlCF}{RECEIVES 1:N\\PERFORMS 1:N}
\newcommand{\ovlDE}{STAY\_LAB 1:N}
\newcommand{\ovlDF}{OFFERS 1:N\\ORDERS 1:N\\ANALYSES 1:N\\UPLOADS 1:N}
\newcommand{\ovlEF}{IN\_DEPT 1:N\\ADMITS 1:N\\ADMITTED\_TO 1:N\\ASSIGNED 1:N\\COVERS 1:N}
\newcommand{\ovcount}{39}

\newcommand{\ovmA}{\textbf{Module 1: Patients and visits}\\[1pt]{\tiny PATIENT, APPOINTMENT, ENCOUNTER, VITAL\textunderscore{}SIGN, CONSULTATION}}
\newcommand{\ovmB}{\textbf{Module 2: Clinical orders}\\[1pt]{\tiny PRESCRIPTION, PRESCRIPTION\textunderscore{}ITEM, MED\textunderscore{}ADMINISTRATION, SURGERY\textunderscore{}REQUEST}}
\newcommand{\ovmC}{\textbf{Module 3: Pharmacy}\\[1pt]{\tiny MEDICINE, MEDICINE\textunderscore{}BATCH, STOCK\textunderscore{}MOVEMENT}}
\newcommand{\ovmD}{\textbf{Module 4: Laboratory}\\[1pt]{\tiny LAB\textunderscore{}TEST, LAB\textunderscore{}ORDER, LAB\textunderscore{}RESULT, LAB\textunderscore{}ATTACHMENT}}
\newcommand{\ovmE}{\textbf{Module 5: Beds and admissions}\\[1pt]{\tiny WARD, BED, ADMISSION, EMERGENCY\textunderscore{}CASE, AMBULANCE, BILL, BILL\textunderscore{}ITEM}}
\newcommand{\ovmF}{\textbf{Module 6: Staff and audit}\\[1pt]{\tiny USER, DOCTOR, DEPARTMENT, DOCTOR\textunderscore{}SCHEDULE, DAILY\textunderscore{}ANALYTICS, WARD\textunderscore{}ASSIGNMENT, AUDIT\textunderscore{}LOG}}

\begin{tikzpicture}[x=1cm,y=0.78cm,
  mod/.style={draw, line width=0.9pt, rounded corners=3pt, fill=#1, align=center, text width=3.7cm, inner sep=4pt, font=\sffamily\scriptsize, execute at begin node={\hyphenpenalty=10000\exhyphenpenalty=10000}},
  link/.style={line width=0.7pt, gray!70!black},
  lab/.style={fill=white, draw=gray!60, rounded corners=1pt, align=center, font=\sffamily\tiny, inner sep=1.5pt}]
  \node[mod=blue!8]   (M1) at (8.5, 0)     {\ovmA};
  \node[mod=orange!10] (M2) at (8.5, -8.3)  {\ovmB};
  \node[mod=yellow!12] (M3) at (8.3, -12.9) {\ovmC};
  \node[mod=violet!8]  (M4) at (14.4, -4.6) {\ovmD};
  \node[mod=green!9]   (M5) at (14.0, -16.4) {\ovmE};
  \node[mod=gray!12]   (M6) at (3.0, -16.4)  {\ovmF};
  \draw[link] (M1) -- node[lab, pos=0.5] {\ovlAB} (M2);
  \draw[link] (M1) -- node[lab, pos=0.42] {\ovlAF} (M6);
  \draw[link] (M1) -- node[lab, pos=0.47] {\ovlAE} (M5);
  \draw[link] (M1) -- node[lab, pos=0.5] {\ovlAD} (M4);
  \draw[link] (M2) -- node[lab, pos=0.55] {\ovlBE} (M5);
  \draw[link] (M2) -- node[lab, pos=0.5] {\ovlBF} (M6);
  \draw[link] (M2) -- node[lab, pos=0.5] {\ovlBC} (M3);
  \draw[link] (M3) -- node[lab, pos=0.5] {\ovlCF} (M6);
  \draw[link] (M4) -- node[lab, pos=0.45] {\ovlDE} (M5);
  \draw[link] (M5) -- node[lab, pos=0.5] {\ovlEF} (M6);
  \draw[link] (M4.east) -- ++(0.5,0) |- node[lab, pos=0.75] {\ovlDF} ($(M6.south)+(0,-1.2)$) -- (M6.south);
\end{tikzpicture}

\caption{Overview: the six modules and every relationship that crosses between two modules (name and cardinality ratio)}\label{fig:overview}
\end{figure}

\begin{figure}[p]\centering
% Generated by docs/report/gen/gen-er.js: diagram 1a (Module 1)
\begin{tikzpicture}[x=1cm,y=1cm]
\node[ent] (PATIENT) at (4.4,-9.2) {PATIENT};
\node[att] (PATIENT-patient_id) at (1.3,-8.2) {\underline{patient\_id}};
\draw (PATIENT-patient_id) -- (PATIENT);
\node[att] (PATIENT-name) at (1.2,-9.4) {name};
\draw (PATIENT-name) -- (PATIENT);
\node[att] (PATIENT-date_of_birth) at (2,-6.9) {date\_of\_birth};
\draw (PATIENT-date_of_birth) -- (PATIENT);
\node[att,dashed] (PATIENT-age) at (4,-6.6) {age};
\draw (PATIENT-age) -- (PATIENT);
\node[att] (PATIENT-gender) at (1.4,-10.6) {gender};
\draw (PATIENT-gender) -- (PATIENT);
\node[att,double,double distance=0.8pt] (PATIENT-allergies) at (1,-7.5) {allergies};
\draw (PATIENT-allergies) -- (PATIENT);
\node[att] (PATIENT-abha_number) at (3.3,-5.7) {abha\_number};
\draw (PATIENT-abha_number) -- (PATIENT);
\node[att] (PATIENT-phone) at (1.6,-11.7) {phone};
\draw (PATIENT-phone) -- (PATIENT);
\node[ent] (APPOINTMENT) at (8.6,-2.4) {APPOINTMENT};
\node[att] (APPOINTMENT-appointment_id) at (5.6,-1) {\underline{appointment\_id}};
\draw (APPOINTMENT-appointment_id) -- (APPOINTMENT);
\node[att] (APPOINTMENT-slot) at (8.6,-0.6) {slot};
\draw (APPOINTMENT-slot) -- (APPOINTMENT);
\node[att] (APPOINTMENT-slot-appointment_date) at (6.9,0.3) {appointment\_date};
\draw (APPOINTMENT-slot-appointment_date) -- (APPOINTMENT-slot);
\node[att] (APPOINTMENT-slot-appointment_time) at (10.3,0.3) {appointment\_time};
\draw (APPOINTMENT-slot-appointment_time) -- (APPOINTMENT-slot);
\node[att] (APPOINTMENT-status) at (11.6,-1) {status};
\draw (APPOINTMENT-status) -- (APPOINTMENT);
\node[att] (APPOINTMENT-reason) at (5.6,-2.4) {reason};
\draw (APPOINTMENT-reason) -- (APPOINTMENT);
\node[ent] (ENCOUNTER) at (13.4,-9.2) {ENCOUNTER};
\node[att] (ENCOUNTER-encounter_id) at (15.5,-6.9) {\underline{encounter\_id}};
\draw (ENCOUNTER-encounter_id) -- (ENCOUNTER);
\node[att] (ENCOUNTER-encounter_type) at (16,-8.1) {encounter\_type};
\draw (ENCOUNTER-encounter_type) -- (ENCOUNTER);
\node[att] (ENCOUNTER-status) at (16,-9.4) {status};
\draw (ENCOUNTER-status) -- (ENCOUNTER);
\node[att] (ENCOUNTER-timeline) at (15.9,-10.7) {timeline};
\draw (ENCOUNTER-timeline) -- (ENCOUNTER);
\node[att] (ENCOUNTER-timeline-arrived_at) at (14.4,-11.3) {arrived\_at};
\draw (ENCOUNTER-timeline-arrived_at) -- (ENCOUNTER-timeline);
\node[att] (ENCOUNTER-timeline-start_timestamp) at (16.1,-11.9) {start\_timestamp};
\draw (ENCOUNTER-timeline-start_timestamp) -- (ENCOUNTER-timeline);
\node[att] (ENCOUNTER-timeline-end_timestamp) at (14.9,-12.6) {end\_timestamp};
\draw (ENCOUNTER-timeline-end_timestamp) -- (ENCOUNTER-timeline);
\node[refent] (DOCTOR) at (14.6,-2.4) {DOCTOR\\[-1pt]{\tiny see Module 6}};
\node[refent] (USER) at (11.2,-6) {USER\\[-1pt]{\tiny see Module 6}};
\node[refent] (PRESCRIPTION) at (8.6,-13.2) {PRESCRIPTION\\[-1pt]{\tiny see Module 2}};
\node[refent] (SURGERY_REQUEST) at (8.6,-15.6) {SURGERY\_REQUEST\\[-1pt]{\tiny see Module 2}};
\node[refent] (LAB_ORDER) at (8.6,-18) {LAB\_ORDER\\[-1pt]{\tiny see Module 4}};
\node[refent] (ADMISSION) at (8.6,-20.6) {ADMISSION\\[-1pt]{\tiny see Module 5}};
\node[refent] (BILL) at (1.4,-15) {BILL\\[-1pt]{\tiny see Module 5}};
\node[refent] (EMERGENCY_CASE) at (1.6,-17.6) {EMERGENCY\_CASE\\[-1pt]{\tiny see Module 5}};
\node[refent] (BED) at (3.4,-22.6) {BED\\[-1pt]{\tiny see Module 5}};
\node[rel] (BOOKS) at (6.3,-5.4) {BOOKS};
\draw[partial] (BOOKS) -- node[card,pos=0.72] {1} (PATIENT);
\draw[total] (BOOKS) -- node[card,pos=0.72] {N} (APPOINTMENT);
\node[rel] (ATTENDS) at (11.6,-2.4) {ATTENDS};
\draw[partial] (ATTENDS) -- node[card,pos=0.72] {1} (DOCTOR);
\draw[total] (ATTENDS) -- node[card,pos=0.72] {N} (APPOINTMENT);
\node[rel] (VISITS) at (8.6,-9.2) {VISITS};
\draw[partial] (VISITS) -- node[card,pos=0.72] {1} (PATIENT);
\draw[total] (VISITS) -- node[card,pos=0.72] {N} (ENCOUNTER);
\node[rel] (CONDUCTS) at (14.6,-5.6) {CONDUCTS};
\draw[partial] (CONDUCTS) -- node[card,pos=0.72] {1} (DOCTOR);
\draw[total] (CONDUCTS) -- node[card,pos=0.72] {N} (ENCOUNTER);
\node[rel] (BECOMES) at (11.6,-4) {BECOMES};
\draw[partial] (BECOMES) -- node[card,pos=0.72] {1} (APPOINTMENT);
\draw[partial] (BECOMES) -- node[card,pos=0.72] {1} (ENCOUNTER);
\node[rel] (OPENS) at (12,-7.4) {OPENS};
\draw[partial] (OPENS) -- node[card,pos=0.72] {1} (USER);
\draw[partial] (OPENS) -- node[card,pos=0.72] {N} (ENCOUNTER);
\node[rel] (STAY_VISIT) at (11.6,-20.6) {STAY\_VISIT};
\draw[partial] (STAY_VISIT) -- node[card,pos=0.72] {1} (ADMISSION);
\draw[partial] (STAY_VISIT) -- node[card,pos=0.72] {N} (ENCOUNTER.-100);
\node[rel] (PRESCRIBED) at (5.6,-13.2) {PRESCRIBED};
\draw[partial] (PRESCRIBED) -- node[card,pos=0.72] {1} (PATIENT.south east);
\draw[total] (PRESCRIBED) -- node[card,pos=0.72] {N} (PRESCRIPTION);
\node[rel] (RX_IN) at (11.6,-13.2) {RX\_IN};
\draw[partial] (RX_IN) -- node[card,pos=0.72] {1} (ENCOUNTER.south west);
\draw[partial] (RX_IN) -- node[card,pos=0.72] {N} (PRESCRIPTION);
\node[rel] (SURGERY_OF) at (5.6,-15.6) {SURGERY\_OF};
\draw[partial] (SURGERY_OF) -- node[card,pos=0.72] {1} (PATIENT.-60);
\draw[total] (SURGERY_OF) -- node[card,pos=0.72] {N} (SURGERY_REQUEST);
\node[rel] (SURGERY_IN) at (11.6,-15.6) {SURGERY\_IN};
\draw[partial] (SURGERY_IN) -- node[card,pos=0.72] {1} (ENCOUNTER.-120);
\draw[partial] (SURGERY_IN) -- node[card,pos=0.72] {N} (SURGERY_REQUEST);
\node[rel] (TESTED) at (5.6,-18) {TESTED};
\draw[partial] (TESTED) -- node[card,pos=0.72] {1} (PATIENT.-70);
\draw[total] (TESTED) -- node[card,pos=0.72] {N} (LAB_ORDER);
\node[rel] (LAB_IN) at (11.6,-18) {LAB\_IN};
\draw[partial] (LAB_IN) -- node[card,pos=0.72] {1} (ENCOUNTER.-110);
\draw[partial] (LAB_IN) -- node[card,pos=0.72] {N} (LAB_ORDER);
\node[rel] (ADMITTED) at (5.6,-20.6) {ADMITTED};
\draw[partial] (ADMITTED) -- node[card,pos=0.72] {1} (PATIENT.-80);
\draw[total] (ADMITTED) -- node[card,pos=0.72] {N} (ADMISSION);
\node[rel] (OCCUPIES) at (6.8,-22.6) {OCCUPIES};
\draw[partial] (OCCUPIES) -- node[card,pos=0.72] {N} (BED);
\draw[total] (OCCUPIES) -- node[card,pos=0.72] {M} (ADMISSION);
\draw[partial] (OCCUPIES.west) -- (4.4,-21.3) -- node[card,pos=0.72] {1} (PATIENT);
\node[att] (OCCUPIES-assigned_at) at (9.2,-23.1) {assigned\_at};
\draw (OCCUPIES-assigned_at) -- (OCCUPIES);
\node[att] (OCCUPIES-released_at) at (11.4,-22.4) {released\_at};
\draw (OCCUPIES-released_at) -- (OCCUPIES);
\node[att] (OCCUPIES-status) at (5,-23.6) {status};
\draw (OCCUPIES-status) -- (OCCUPIES);
\node[rel] (ARRIVES) at (3.4,-16.2) {ARRIVES};
\draw[partial] (ARRIVES) -- node[card,pos=0.72] {1} (PATIENT.-115);
\draw[partial] (ARRIVES) -- node[card,pos=0.72] {N} (EMERGENCY_CASE);
\node[rel] (BILLED) at (3.4,-13.6) {BILLED};
\draw[partial] (BILLED) -- node[card,pos=0.72] {1} (PATIENT.south west);
\draw[total] (BILLED) -- node[card,pos=0.72] {N} (BILL);
\end{tikzpicture}
\minilegend
\caption{Module 1, Patients and visits (page 1 of 2): patients, appointments and encounters}\label{fig:m1a}
\end{figure}

\begin{figure}[p]\centering
% Generated by docs/report/gen/gen-er.js: diagram 1b (Module 1)
\begin{tikzpicture}[x=1cm,y=0.72cm]
\node[ent] (VITAL_SIGN) at (3.2,-9.2) {VITAL\_SIGN};
\node[att] (VITAL_SIGN-vital_id) at (1.1,-11) {\underline{vital\_id}};
\draw (VITAL_SIGN-vital_id) -- (VITAL_SIGN);
\node[att] (VITAL_SIGN-recorded_at_utc) at (1.4,-12.1) {recorded\_at\_utc};
\draw (VITAL_SIGN-recorded_at_utc) -- (VITAL_SIGN);
\node[att] (VITAL_SIGN-blood_pressure) at (3,-14) {blood\_pressure};
\draw (VITAL_SIGN-blood_pressure) -- (VITAL_SIGN);
\node[att] (VITAL_SIGN-blood_pressure-bp_systolic) at (1.8,-15.2) {bp\_systolic};
\draw (VITAL_SIGN-blood_pressure-bp_systolic) -- (VITAL_SIGN-blood_pressure);
\node[att] (VITAL_SIGN-blood_pressure-bp_diastolic) at (4.3,-15.2) {bp\_diastolic};
\draw (VITAL_SIGN-blood_pressure-bp_diastolic) -- (VITAL_SIGN-blood_pressure);
\node[att] (VITAL_SIGN-pulse_bpm) at (3.6,-12.4) {pulse\_bpm};
\draw (VITAL_SIGN-pulse_bpm) -- (VITAL_SIGN);
\node[att] (VITAL_SIGN-temperature_c) at (1,-9.9) {temperature\_c};
\draw (VITAL_SIGN-temperature_c) -- (VITAL_SIGN);
\node[att] (VITAL_SIGN-spo2_pct) at (1,-7.9) {spo2\_pct};
\draw (VITAL_SIGN-spo2_pct) -- (VITAL_SIGN);
\node[ent] (CONSULTATION) at (13.8,-9.2) {CONSULTATION};
\node[att] (CONSULTATION-consultation_id) at (16.2,-10.4) {\underline{consultation\_id}};
\draw (CONSULTATION-consultation_id) -- (CONSULTATION);
\node[att] (CONSULTATION-consultation_time) at (16,-11.5) {consultation\_time};
\draw (CONSULTATION-consultation_time) -- (CONSULTATION);
\node[att] (CONSULTATION-symptoms) at (16.2,-9.2) {symptoms};
\draw (CONSULTATION-symptoms) -- (CONSULTATION);
\node[att] (CONSULTATION-diagnosis) at (16.2,-12.6) {diagnosis};
\draw (CONSULTATION-diagnosis) -- (CONSULTATION);
\node[att] (CONSULTATION-plan) at (15.9,-13.7) {plan};
\draw (CONSULTATION-plan) -- (CONSULTATION);
\node[refent] (PATIENT) at (8.5,-4.2) {PATIENT\\[-1pt]{\tiny see Module 1, page 1}};
\node[refent] (ENCOUNTER) at (8.5,-9.2) {ENCOUNTER\\[-1pt]{\tiny see Module 1, page 1}};
\node[refent] (APPOINTMENT) at (13.8,-1.4) {APPOINTMENT\\[-1pt]{\tiny see Module 1, page 1}};
\node[refent] (ADMISSION) at (8.5,-15) {ADMISSION\\[-1pt]{\tiny see Module 5}};
\node[refent] (USER) at (3.2,-1.4) {USER\\[-1pt]{\tiny see Module 6}};
\node[refent] (DOCTOR) at (15.299999999999999,-4.6) {DOCTOR\\[-1pt]{\tiny see Module 6}};
\node[refent] (PRESCRIPTION) at (11,-18.4) {PRESCRIPTION\\[-1pt]{\tiny see Module 2}};
\node[refent] (LAB_ORDER) at (15.2,-18.4) {LAB\_ORDER\\[-1pt]{\tiny see Module 4}};
\node[rel] (MEASURED) at (5.8,-6.3) {MEASURED};
\draw[partial] (MEASURED) -- node[card,pos=0.72] {1} (PATIENT);
\draw[total] (MEASURED) -- node[card,pos=0.72] {N} (VITAL_SIGN);
\node[rel] (TAKEN_IN) at (5.8,-9.2) {TAKEN\_IN};
\draw[partial] (TAKEN_IN) -- node[card,pos=0.72] {1} (ENCOUNTER);
\draw[partial] (TAKEN_IN) -- node[card,pos=0.72] {N} (VITAL_SIGN);
\node[rel] (TAKEN_ON) at (5.8,-13.2) {TAKEN\_ON};
\draw[partial] (TAKEN_ON) -- node[card,pos=0.72] {1} (ADMISSION);
\draw[partial] (TAKEN_ON) -- node[card,pos=0.72] {N} (VITAL_SIGN);
\node[rel] (RECORDS) at (3.2,-5) {RECORDS};
\draw[partial] (RECORDS) -- node[card,pos=0.72] {1} (USER);
\draw[total] (RECORDS) -- node[card,pos=0.72] {N} (VITAL_SIGN);
\node[rel] (NOTE_OF) at (11.2,-6.3) {NOTE\_OF};
\draw[partial] (NOTE_OF) -- node[card,pos=0.72] {1} (PATIENT);
\draw[total] (NOTE_OF) -- node[card,pos=0.72] {N} (CONSULTATION);
\node[rel] (WRITES) at (16.1,-6.4) {WRITES};
\draw[partial] (WRITES) -- node[card,pos=0.72] {1} (DOCTOR);
\draw[total] (WRITES) -- node[card,pos=0.72] {N} (CONSULTATION);
\node[rel] (NOTE_IN) at (11.2,-9.2) {NOTE\_IN};
\draw[partial] (NOTE_IN) -- node[card,pos=0.72] {1} (ENCOUNTER);
\draw[partial] (NOTE_IN) -- node[card,pos=0.72] {N} (CONSULTATION);
\node[rel] (NOTE_FOR) at (13.8,-4.2) {NOTE\_FOR};
\draw[partial] (NOTE_FOR) -- node[card,pos=0.72] {1} (APPOINTMENT);
\draw[partial] (NOTE_FOR) -- node[card,pos=0.72] {N} (CONSULTATION);
\node[rel] (STAY_NOTE) at (11.2,-13.2) {STAY\_NOTE};
\draw[partial] (STAY_NOTE) -- node[card,pos=0.72] {1} (ADMISSION);
\draw[partial] (STAY_NOTE) -- node[card,pos=0.72] {N} (CONSULTATION);
\node[rel] (RX_FROM) at (12,-15.6) {RX\_FROM};
\draw[partial] (RX_FROM) -- node[card,pos=0.72] {1} (CONSULTATION);
\draw[partial] (RX_FROM) -- node[card,pos=0.72] {N} (PRESCRIPTION);
\node[rel] (LAB_FROM) at (15.2,-15.6) {LAB\_FROM};
\draw[partial] (LAB_FROM) -- node[card,pos=0.72] {1} (CONSULTATION);
\draw[partial] (LAB_FROM) -- node[card,pos=0.72] {N} (LAB_ORDER);
\end{tikzpicture}
\minilegend
\caption{Module 1, Patients and visits (page 2 of 2): vital signs and consultations}\label{fig:m1b}
\vspace{5mm}
% Generated by docs/report/gen/gen-er.js: diagram 3 (Module 3)
\begin{tikzpicture}[x=1cm,y=0.7cm]
\node[ent] (MEDICINE) at (6,-4) {MEDICINE};
\node[att] (MEDICINE-medicine_id) at (3.4,-2.6) {\underline{medicine\_id}};
\draw (MEDICINE-medicine_id) -- (MEDICINE);
\node[att] (MEDICINE-name) at (4.4,-1.4) {name};
\draw (MEDICINE-name) -- (MEDICINE);
\node[att] (MEDICINE-category) at (6.4,-1) {category};
\draw (MEDICINE-category) -- (MEDICINE);
\node[att] (MEDICINE-unit_price) at (8.6,-1.6) {unit\_price};
\draw (MEDICINE-unit_price) -- (MEDICINE);
\node[att] (MEDICINE-reorder_level) at (9.6,-2.8) {reorder\_level};
\draw (MEDICINE-reorder_level) -- (MEDICINE);
\node[att,dashed] (MEDICINE-stock_quantity) at (3.3,-5) {stock\_quantity};
\draw (MEDICINE-stock_quantity) -- (MEDICINE);
\node[weakent] (MEDICINE_BATCH) at (6,-11) {MEDICINE\_BATCH};
\node[att] (MEDICINE_BATCH-batch_number) at (3,-12.4) {\dashuline{batch\_number}};
\draw (MEDICINE_BATCH-batch_number) -- (MEDICINE_BATCH);
\node[att] (MEDICINE_BATCH-quantity) at (4,-13.6) {quantity};
\draw (MEDICINE_BATCH-quantity) -- (MEDICINE_BATCH);
\node[att] (MEDICINE_BATCH-expiry_date) at (6,-14) {expiry\_date};
\draw (MEDICINE_BATCH-expiry_date) -- (MEDICINE_BATCH);
\node[att] (MEDICINE_BATCH-received_date) at (3.3,-10) {received\_date};
\draw (MEDICINE_BATCH-received_date) -- (MEDICINE_BATCH);
\node[ent] (STOCK_MOVEMENT) at (13.2,-11) {STOCK\_MOVEMENT};
\node[att] (STOCK_MOVEMENT-movement_id) at (15.6,-9.4) {\underline{movement\_id}};
\draw (STOCK_MOVEMENT-movement_id) -- (STOCK_MOVEMENT);
\node[att] (STOCK_MOVEMENT-movement_type) at (16,-10.6) {movement\_type};
\draw (STOCK_MOVEMENT-movement_type) -- (STOCK_MOVEMENT);
\node[att] (STOCK_MOVEMENT-quantity) at (16,-11.8) {quantity};
\draw (STOCK_MOVEMENT-quantity) -- (STOCK_MOVEMENT);
\node[att] (STOCK_MOVEMENT-reason) at (15.6,-13) {reason};
\draw (STOCK_MOVEMENT-reason) -- (STOCK_MOVEMENT);
\node[att] (STOCK_MOVEMENT-created_at) at (15.4,-8.3) {created\_at};
\draw (STOCK_MOVEMENT-created_at) -- (STOCK_MOVEMENT);
\node[refent] (PRESCRIPTION_ITEM) at (1.4,-7.6) {PRESCRIPTION\_ITEM\\[-1pt]{\tiny see Module 2}};
\node[refent] (USER) at (11,-16) {USER\\[-1pt]{\tiny see Module 6}};
\node[rel] (OF_DRUG) at (1.6,-4) {OF\_DRUG};
\draw[partial] (OF_DRUG) -- node[card,pos=0.72] {1} (MEDICINE.west);
\draw[total] (OF_DRUG) -- node[card,pos=0.72] {N} (PRESCRIPTION_ITEM);
\node[rel] (FILLED_FROM) at (1.6,-11) {FILLED\_FROM};
\draw[partial] (FILLED_FROM) -- node[card,pos=0.72] {M} (PRESCRIPTION_ITEM);
\draw[partial] (FILLED_FROM) -- node[card,pos=0.72] {N} (MEDICINE_BATCH.west);
\node[att] (FILLED_FROM-quantity) at (1.4,-12.6) {quantity};
\draw (FILLED_FROM-quantity) -- (FILLED_FROM);
\node[idrel] (HAS_BATCH) at (6,-7.5) {HAS\_BATCH};
\draw[partial] (HAS_BATCH) -- node[card,pos=0.72] {1} (MEDICINE);
\draw[total] (HAS_BATCH) -- node[card,pos=0.72] {N} (MEDICINE_BATCH);
\node[rel] (RECEIVES) at (8,-14.2) {RECEIVES};
\draw[partial] (RECEIVES) -- node[card,pos=0.72] {1} (USER);
\draw[partial] (RECEIVES) -- node[card,pos=0.72] {N} (MEDICINE_BATCH.south east);
\node[rel] (MOVES) at (11.2,-6.4) {MOVES};
\draw[partial] (MOVES) -- node[card,pos=0.72] {1} (MEDICINE);
\draw[total] (MOVES) -- node[card,pos=0.72] {N} (STOCK_MOVEMENT);
\node[rel] (FROM_BATCH) at (9.6,-11) {FROM\_BATCH};
\draw[partial] (FROM_BATCH) -- node[card,pos=0.72] {1} (MEDICINE_BATCH);
\draw[partial] (FROM_BATCH) -- node[card,pos=0.72] {N} (STOCK_MOVEMENT);
\node[rel] (PERFORMS) at (13.2,-14) {PERFORMS};
\draw[partial] (PERFORMS) -- node[card,pos=0.72] {1} (USER);
\draw[total] (PERFORMS) -- node[card,pos=0.72] {N} (STOCK_MOVEMENT);
\end{tikzpicture}
\minilegend
\caption{Module 3, Pharmacy: medicines, batches and stock movements}\label{fig:m3}
\end{figure}

\begin{figure}[p]\centering
% Generated by docs/report/gen/gen-er.js: diagram 2 (Module 2)
\begin{tikzpicture}[x=1cm,y=1cm]
\node[ent] (PRESCRIPTION) at (3.6,-7.4) {PRESCRIPTION};
\node[att] (PRESCRIPTION-prescription_id) at (1.1,-6.2) {\underline{prescription\_id}};
\draw (PRESCRIPTION-prescription_id) -- (PRESCRIPTION);
\node[att] (PRESCRIPTION-prescription_date) at (1.2,-8.6) {prescription\_date};
\draw (PRESCRIPTION-prescription_date) -- (PRESCRIPTION);
\node[att] (PRESCRIPTION-status) at (1.3,-9.7) {status};
\draw (PRESCRIPTION-status) -- (PRESCRIPTION);
\node[att] (PRESCRIPTION-notes) at (1,-7.4) {notes};
\draw (PRESCRIPTION-notes) -- (PRESCRIPTION);
\node[ent] (PRESCRIPTION_ITEM) at (3.6,-13.6) {PRESCRIPTION\_ITEM};
\node[att] (PRESCRIPTION_ITEM-item_id) at (1,-12.4) {\underline{item\_id}};
\draw (PRESCRIPTION_ITEM-item_id) -- (PRESCRIPTION_ITEM);
\node[att] (PRESCRIPTION_ITEM-dosage) at (1,-13.6) {dosage};
\draw (PRESCRIPTION_ITEM-dosage) -- (PRESCRIPTION_ITEM);
\node[att] (PRESCRIPTION_ITEM-frequency_code) at (1.3,-14.8) {frequency\_code};
\draw (PRESCRIPTION_ITEM-frequency_code) -- (PRESCRIPTION_ITEM);
\node[att] (PRESCRIPTION_ITEM-quantity) at (1.4,-15.9) {quantity};
\draw (PRESCRIPTION_ITEM-quantity) -- (PRESCRIPTION_ITEM);
\node[att] (PRESCRIPTION_ITEM-dispensed) at (6,-15.2) {dispensed};
\draw (PRESCRIPTION_ITEM-dispensed) -- (PRESCRIPTION_ITEM);
\node[att] (PRESCRIPTION_ITEM-dispensed_at) at (6.4,-16.2) {dispensed\_at};
\draw (PRESCRIPTION_ITEM-dispensed_at) -- (PRESCRIPTION_ITEM);
\node[ent] (MED_ADMINISTRATION) at (13.2,-13.6) {MED\_ADMINISTRATION};
\node[att] (MED_ADMINISTRATION-administration_id) at (15.6,-12.2) {\underline{administration\_id}};
\draw (MED_ADMINISTRATION-administration_id) -- (MED_ADMINISTRATION);
\node[att] (MED_ADMINISTRATION-scheduled_at_utc) at (15.9,-13.4) {scheduled\_at\_utc};
\draw (MED_ADMINISTRATION-scheduled_at_utc) -- (MED_ADMINISTRATION);
\node[att] (MED_ADMINISTRATION-status) at (15.9,-14.6) {status};
\draw (MED_ADMINISTRATION-status) -- (MED_ADMINISTRATION);
\node[att] (MED_ADMINISTRATION-dose_given) at (15.5,-15.7) {dose\_given};
\draw (MED_ADMINISTRATION-dose_given) -- (MED_ADMINISTRATION);
\node[att] (MED_ADMINISTRATION-route_given) at (11,-15.4) {route\_given};
\draw (MED_ADMINISTRATION-route_given) -- (MED_ADMINISTRATION);
\node[ent] (SURGERY_REQUEST) at (14.6,-4.6) {SURGERY\_REQUEST};
\node[att] (SURGERY_REQUEST-request_id) at (13.4,-1.2) {\underline{request\_id}};
\draw (SURGERY_REQUEST-request_id) -- (SURGERY_REQUEST);
\node[att] (SURGERY_REQUEST-procedure_name) at (15.7,-1.9) {procedure\_name};
\draw (SURGERY_REQUEST-procedure_name) -- (SURGERY_REQUEST);
\node[att] (SURGERY_REQUEST-priority) at (16,-6.4) {priority};
\draw (SURGERY_REQUEST-priority) -- (SURGERY_REQUEST);
\node[att] (SURGERY_REQUEST-status) at (14.6,-7.4) {status};
\draw (SURGERY_REQUEST-status) -- (SURGERY_REQUEST);
\node[refent] (PATIENT) at (9,-2) {PATIENT\\[-1pt]{\tiny see Module 1}};
\node[refent] (DOCTOR) at (9,-4.6) {DOCTOR\\[-1pt]{\tiny see Module 6}};
\node[refent] (ENCOUNTER) at (9,-7.2) {ENCOUNTER\\[-1pt]{\tiny see Module 1}};
\node[refent] (CONSULTATION) at (2,-1.6) {CONSULTATION\\[-1pt]{\tiny see Module 1}};
\node[refent] (ADMISSION) at (9,-10.6) {ADMISSION\\[-1pt]{\tiny see Module 5}};
\node[refent] (MEDICINE) at (2.6,-21) {MEDICINE\\[-1pt]{\tiny see Module 3}};
\node[refent] (MEDICINE_BATCH) at (8.2,-21) {MEDICINE\_BATCH\\[-1pt]{\tiny see Module 3}};
\node[refent] (USER) at (13.2,-21) {USER\\[-1pt]{\tiny see Module 6}};
\node[rel] (PRESCRIBED) at (6.2,-2) {PRESCRIBED};
\draw[partial] (PRESCRIBED) -- node[card,pos=0.72] {1} (PATIENT);
\draw[total] (PRESCRIBED) -- node[card,pos=0.72] {N} (PRESCRIPTION.75);
\node[rel] (PRESCRIBES) at (6.2,-4.6) {PRESCRIBES};
\draw[partial] (PRESCRIBES) -- node[card,pos=0.72] {1} (DOCTOR);
\draw[total] (PRESCRIBES) -- node[card,pos=0.72] {N} (PRESCRIPTION.north east);
\node[rel] (RX_IN) at (6.2,-7.2) {RX\_IN};
\draw[partial] (RX_IN) -- node[card,pos=0.72] {1} (ENCOUNTER);
\draw[partial] (RX_IN) -- node[card,pos=0.72] {N} (PRESCRIPTION);
\node[rel] (RX_FROM) at (2.4,-4.3) {RX\_FROM};
\draw[partial] (RX_FROM) -- node[card,pos=0.72] {1} (CONSULTATION);
\draw[partial] (RX_FROM) -- node[card,pos=0.72] {N} (PRESCRIPTION.110);
\node[rel] (STAY_RX) at (6.2,-9.8) {STAY\_RX};
\draw[partial] (STAY_RX) -- node[card,pos=0.72] {1} (ADMISSION);
\draw[partial] (STAY_RX) -- node[card,pos=0.72] {N} (PRESCRIPTION.south east);
\node[rel] (LISTS) at (3.6,-10.6) {LISTS};
\draw[partial] (LISTS) -- node[card,pos=0.72] {1} (PRESCRIPTION);
\draw[total] (LISTS) -- node[card,pos=0.72] {N} (PRESCRIPTION_ITEM);
\node[rel] (OF_DRUG) at (2.6,-18.6) {OF\_DRUG};
\draw[partial] (OF_DRUG) -- node[card,pos=0.72] {1} (MEDICINE);
\draw[total] (OF_DRUG) -- node[card,pos=0.72] {N} (PRESCRIPTION_ITEM.-110);
\node[rel] (FILLED_FROM) at (6.8,-18.2) {FILLED\_FROM};
\draw[partial] (FILLED_FROM) -- node[card,pos=0.72] {M} (PRESCRIPTION_ITEM.south east);
\draw[partial] (FILLED_FROM) -- node[card,pos=0.72] {N} (MEDICINE_BATCH);
\node[att] (FILLED_FROM-quantity) at (9,-17.4) {quantity};
\draw (FILLED_FROM-quantity) -- (FILLED_FROM);
\node[rel] (GIVEN_AS) at (8.4,-13.6) {GIVEN\_AS};
\draw[partial] (GIVEN_AS) -- node[card,pos=0.72] {1} (PRESCRIPTION_ITEM);
\draw[total] (GIVEN_AS) -- node[card,pos=0.72] {N} (MED_ADMINISTRATION);
\node[rel] (STAY_DOSE) at (11.6,-10.6) {STAY\_DOSE};
\draw[partial] (STAY_DOSE) -- node[card,pos=0.72] {1} (ADMISSION);
\draw[partial] (STAY_DOSE) -- node[card,pos=0.72] {N} (MED_ADMINISTRATION.north west);
\node[rel] (GIVEN_BY) at (13.2,-18.3) {GIVEN\_BY};
\draw[partial] (GIVEN_BY) -- node[card,pos=0.72] {1} (USER);
\draw[partial] (GIVEN_BY) -- node[card,pos=0.72] {N} (MED_ADMINISTRATION);
\node[rel] (SURGERY_OF) at (11.8,-2) {SURGERY\_OF};
\draw[partial] (SURGERY_OF) -- node[card,pos=0.72] {1} (PATIENT);
\draw[total] (SURGERY_OF) -- node[card,pos=0.72] {N} (SURGERY_REQUEST.north west);
\node[rel] (REQUESTS) at (11.8,-4.6) {REQUESTS};
\draw[partial] (REQUESTS) -- node[card,pos=0.72] {1} (DOCTOR);
\draw[total] (REQUESTS) -- node[card,pos=0.72] {N} (SURGERY_REQUEST);
\node[rel] (SURGERY_IN) at (11.8,-7.2) {SURGERY\_IN};
\draw[partial] (SURGERY_IN) -- node[card,pos=0.72] {1} (ENCOUNTER);
\draw[partial] (SURGERY_IN) -- node[card,pos=0.72] {N} (SURGERY_REQUEST.south west);
\end{tikzpicture}
\minilegend
\caption{Module 2, Clinical orders: prescriptions, prescription lines, medication administration and surgery requests}\label{fig:m2}
\end{figure}

\begin{figure}[p]\centering
% Generated by docs/report/gen/gen-er.js: diagram 4 (Module 4)
\begin{tikzpicture}[x=1cm,y=1cm]
\node[ent] (LAB_TEST) at (2.8,-8.6) {LAB\_TEST};
\node[att] (LAB_TEST-test_id) at (1.1,-10.6) {\underline{test\_id}};
\draw (LAB_TEST-test_id) -- (LAB_TEST);
\node[att] (LAB_TEST-name) at (2.6,-11.4) {name};
\draw (LAB_TEST-name) -- (LAB_TEST);
\node[att] (LAB_TEST-cost) at (4.3,-10.9) {cost};
\draw (LAB_TEST-cost) -- (LAB_TEST);
\node[att] (LAB_TEST-unit) at (1,-6.8) {unit};
\draw (LAB_TEST-unit) -- (LAB_TEST);
\node[att] (LAB_TEST-reference_range) at (3,-13) {reference\_range};
\draw (LAB_TEST-reference_range) -- (LAB_TEST);
\node[att] (LAB_TEST-reference_range-reference_low) at (1.6,-14.2) {reference\_low};
\draw (LAB_TEST-reference_range-reference_low) -- (LAB_TEST-reference_range);
\node[att] (LAB_TEST-reference_range-reference_high) at (4.4,-14.2) {reference\_high};
\draw (LAB_TEST-reference_range-reference_high) -- (LAB_TEST-reference_range);
\node[ent] (LAB_ORDER) at (8.6,-8.6) {LAB\_ORDER};
\node[att] (LAB_ORDER-order_id) at (6.3,-6.4) {\underline{order\_id}};
\draw (LAB_ORDER-order_id) -- (LAB_ORDER);
\node[att] (LAB_ORDER-order_date) at (7.4,-5.3) {order\_date};
\draw (LAB_ORDER-order_date) -- (LAB_ORDER);
\node[att] (LAB_ORDER-status) at (9.2,-5.2) {status};
\draw (LAB_ORDER-status) -- (LAB_ORDER);
\node[att] (LAB_ORDER-priority) at (6.1,-7.5) {priority};
\draw (LAB_ORDER-priority) -- (LAB_ORDER);
\node[ent] (LAB_RESULT) at (8.6,-15) {LAB\_RESULT};
\node[att] (LAB_RESULT-result_id) at (5.8,-14) {\underline{result\_id}};
\draw (LAB_RESULT-result_id) -- (LAB_RESULT);
\node[att] (LAB_RESULT-result_value) at (5.6,-15.2) {result\_value};
\draw (LAB_RESULT-result_value) -- (LAB_RESULT);
\node[att] (LAB_RESULT-numeric_value) at (5.9,-16.4) {numeric\_value};
\draw (LAB_RESULT-numeric_value) -- (LAB_RESULT);
\node[att] (LAB_RESULT-interpretation) at (10.9,-13.9) {interpretation};
\draw (LAB_RESULT-interpretation) -- (LAB_RESULT);
\node[att] (LAB_RESULT-result_date) at (11.2,-15) {result\_date};
\draw (LAB_RESULT-result_date) -- (LAB_RESULT);
\node[ent] (LAB_ATTACHMENT) at (4.6,-20.4) {LAB\_ATTACHMENT};
\node[att] (LAB_ATTACHMENT-attachment_id) at (1.6,-19.4) {\underline{attachment\_id}};
\draw (LAB_ATTACHMENT-attachment_id) -- (LAB_ATTACHMENT);
\node[att] (LAB_ATTACHMENT-original_filename) at (1.8,-21.5) {original\_filename};
\draw (LAB_ATTACHMENT-original_filename) -- (LAB_ATTACHMENT);
\node[att] (LAB_ATTACHMENT-mime_type) at (4.6,-22.2) {mime\_type};
\draw (LAB_ATTACHMENT-mime_type) -- (LAB_ATTACHMENT);
\node[att] (LAB_ATTACHMENT-file_size) at (2,-20.4) {file\_size};
\draw (LAB_ATTACHMENT-file_size) -- (LAB_ATTACHMENT);
\node[refent] (DEPARTMENT) at (2.8,-2.2) {DEPARTMENT\\[-1pt]{\tiny see Module 6}};
\node[refent] (PATIENT) at (15.3,-2.6) {PATIENT\\[-1pt]{\tiny see Module 1}};
\node[refent] (DOCTOR) at (15.3,-5.6) {DOCTOR\\[-1pt]{\tiny see Module 6}};
\node[refent] (ENCOUNTER) at (15.3,-8.6) {ENCOUNTER\\[-1pt]{\tiny see Module 1}};
\node[refent] (CONSULTATION) at (15.3,-11.6) {CONSULTATION\\[-1pt]{\tiny see Module 1}};
\node[refent] (ADMISSION) at (15.3,-14.6) {ADMISSION\\[-1pt]{\tiny see Module 5}};
\node[refent] (USER) at (12.6,-20.4) {USER\\[-1pt]{\tiny see Module 6}};
\node[rel] (OF_TEST) at (5.7,-8.6) {OF\_TEST};
\draw[partial] (OF_TEST) -- node[card,pos=0.72] {1} (LAB_TEST);
\draw[total] (OF_TEST) -- node[card,pos=0.72] {N} (LAB_ORDER);
\node[rel] (OFFERS) at (2.8,-5.1) {OFFERS};
\draw[partial] (OFFERS) -- node[card,pos=0.72] {1} (DEPARTMENT);
\draw[partial] (OFFERS) -- node[card,pos=0.72] {N} (LAB_TEST);
\node[rel] (TESTED) at (12.3,-2.6) {TESTED};
\draw[partial] (TESTED) -- node[card,pos=0.72] {1} (PATIENT);
\draw[total] (TESTED) -- node[card,pos=0.72] {N} (LAB_ORDER.north east);
\node[rel] (ORDERS) at (12.3,-5.6) {ORDERS};
\draw[partial] (ORDERS) -- node[card,pos=0.72] {1} (DOCTOR);
\draw[total] (ORDERS) -- node[card,pos=0.72] {N} (LAB_ORDER.20);
\node[rel] (LAB_IN) at (12.3,-8.6) {LAB\_IN};
\draw[partial] (LAB_IN) -- node[card,pos=0.72] {1} (ENCOUNTER);
\draw[partial] (LAB_IN) -- node[card,pos=0.72] {N} (LAB_ORDER);
\node[rel] (LAB_FROM) at (12.3,-11.6) {LAB\_FROM};
\draw[partial] (LAB_FROM) -- node[card,pos=0.72] {1} (CONSULTATION);
\draw[partial] (LAB_FROM) -- node[card,pos=0.72] {N} (LAB_ORDER.-20);
\node[rel] (STAY_LAB) at (12.3,-14.6) {STAY\_LAB};
\draw[partial] (STAY_LAB) -- node[card,pos=0.72] {1} (ADMISSION);
\draw[partial] (STAY_LAB) -- node[card,pos=0.8] {N} (LAB_ORDER.south east);
\node[rel] (YIELDS) at (8.6,-11.8) {YIELDS};
\draw[partial] (YIELDS) -- node[card,pos=0.72] {1} (LAB_ORDER);
\draw[total] (YIELDS) -- node[card,pos=0.72] {1} (LAB_RESULT);
\node[rel] (ANALYSES) at (11.6,-17.8) {ANALYSES};
\draw[partial] (ANALYSES) -- node[card,pos=0.72] {1} (USER);
\draw[partial] (ANALYSES) -- node[card,pos=0.72] {N} (LAB_RESULT.south east);
\node[rel] (ATTACHED) at (6.6,-17.8) {ATTACHED};
\draw[partial] (ATTACHED) -- node[card,pos=0.72] {1} (LAB_RESULT.south west);
\draw[total] (ATTACHED) -- node[card,pos=0.72] {N} (LAB_ATTACHMENT);
\node[rel] (UPLOADS) at (8.6,-20.4) {UPLOADS};
\draw[partial] (UPLOADS) -- node[card,pos=0.72] {1} (USER);
\draw[total] (UPLOADS) -- node[card,pos=0.72] {N} (LAB_ATTACHMENT);
\end{tikzpicture}
\minilegend
\caption{Module 4, Laboratory: tests, orders, results and report files}\label{fig:m4}
\end{figure}

\begin{figure}[p]\centering
% Generated by docs/report/gen/gen-er.js: diagram 5a (Module 5)
\begin{tikzpicture}[x=1cm,y=1cm]
\node[ent] (WARD) at (3.4,-4.4) {WARD};
\node[att] (WARD-ward_id) at (1,-3) {\underline{ward\_id}};
\draw (WARD-ward_id) -- (WARD);
\node[att] (WARD-name) at (1,-4.2) {name};
\draw (WARD-name) -- (WARD);
\node[att] (WARD-floor) at (1,-5.4) {floor};
\draw (WARD-floor) -- (WARD);
\node[att] (WARD-capacity) at (2.6,-6.4) {capacity};
\draw (WARD-capacity) -- (WARD);
\node[att] (WARD-ward_type) at (4.6,-6.4) {ward\_type};
\draw (WARD-ward_type) -- (WARD);
\node[weakent] (BED) at (3.4,-11) {BED};
\node[att] (BED-bed_number) at (1,-9.6) {\dashuline{bed\_number}};
\draw (BED-bed_number) -- (BED);
\node[att] (BED-bed_type) at (1,-10.8) {bed\_type};
\draw (BED-bed_type) -- (BED);
\node[att] (BED-status) at (1,-12) {status};
\draw (BED-status) -- (BED);
\node[att] (BED-has_ventilator) at (1.3,-13.2) {has\_ventilator};
\draw (BED-has_ventilator) -- (BED);
\node[ent] (ADMISSION) at (11.2,-11) {ADMISSION};
\node[att] (ADMISSION-admission_id) at (12.8,-8.4) {\underline{admission\_id}};
\draw (ADMISSION-admission_id) -- (ADMISSION);
\node[att] (ADMISSION-admission_date) at (15.9,-9.4) {admission\_date};
\draw (ADMISSION-admission_date) -- (ADMISSION);
\node[att] (ADMISSION-discharge_date) at (16,-10.5) {discharge\_date};
\draw (ADMISSION-discharge_date) -- (ADMISSION);
\node[att] (ADMISSION-status) at (15.8,-11.6) {status};
\draw (ADMISSION-status) -- (ADMISSION);
\node[att] (ADMISSION-diagnosis) at (14.6,-8.2) {diagnosis};
\draw (ADMISSION-diagnosis) -- (ADMISSION);
\node[ent] (EMERGENCY_CASE) at (11.2,-18.6) {EMERGENCY\_CASE};
\node[att] (EMERGENCY_CASE-emergency_id) at (15.7,-17.9) {\underline{emergency\_id}};
\draw (EMERGENCY_CASE-emergency_id) -- (EMERGENCY_CASE);
\node[att] (EMERGENCY_CASE-severity) at (15.9,-19) {severity};
\draw (EMERGENCY_CASE-severity) -- (EMERGENCY_CASE);
\node[att] (EMERGENCY_CASE-arrival_time) at (15.6,-20.1) {arrival\_time};
\draw (EMERGENCY_CASE-arrival_time) -- (EMERGENCY_CASE);
\node[att] (EMERGENCY_CASE-status) at (14.4,-21.1) {status};
\draw (EMERGENCY_CASE-status) -- (EMERGENCY_CASE);
\node[att] (EMERGENCY_CASE-requirements) at (11.2,-21.2) {requirements};
\draw (EMERGENCY_CASE-requirements) -- (EMERGENCY_CASE);
\node[att] (EMERGENCY_CASE-requirements-required_bed_type) at (8.6,-22.2) {required\_bed\_type};
\draw (EMERGENCY_CASE-requirements-required_bed_type) -- (EMERGENCY_CASE-requirements);
\node[att] (EMERGENCY_CASE-requirements-required_specialization) at (11.4,-22.8) {required\_specialization};
\draw (EMERGENCY_CASE-requirements-required_specialization) -- (EMERGENCY_CASE-requirements);
\node[att] (EMERGENCY_CASE-requirements-ventilator_required) at (14.4,-22.2) {ventilator\_required};
\draw (EMERGENCY_CASE-requirements-ventilator_required) -- (EMERGENCY_CASE-requirements);
\node[ent] (AMBULANCE) at (3.4,-21.2) {AMBULANCE};
\node[att] (AMBULANCE-ambulance_id) at (1,-19.8) {\underline{ambulance\_id}};
\draw (AMBULANCE-ambulance_id) -- (AMBULANCE);
\node[att] (AMBULANCE-vehicle_number) at (1.2,-22.6) {vehicle\_number};
\draw (AMBULANCE-vehicle_number) -- (AMBULANCE);
\node[att] (AMBULANCE-driver) at (5.4,-23) {driver};
\draw (AMBULANCE-driver) -- (AMBULANCE);
\node[att] (AMBULANCE-driver-driver_name) at (4.4,-24.1) {driver\_name};
\draw (AMBULANCE-driver-driver_name) -- (AMBULANCE-driver);
\node[att] (AMBULANCE-driver-driver_phone) at (6.6,-24.1) {driver\_phone};
\draw (AMBULANCE-driver-driver_phone) -- (AMBULANCE-driver);
\node[att] (AMBULANCE-status) at (3.4,-23.2) {status};
\draw (AMBULANCE-status) -- (AMBULANCE);
\node[refent] (DEPARTMENT) at (11.2,-4.4) {DEPARTMENT\\[-1pt]{\tiny see Module 6}};
\node[refent] (WARD_ASSIGNMENT) at (3.4,-0.6) {WARD\_ASSIGNMENT\\[-1pt]{\tiny see Module 6}};
\node[refent] (PATIENT) at (7.3,-15.9) {PATIENT\\[-1pt]{\tiny see Module 1}};
\node[refent] (DOCTOR) at (15.4,-14.8) {DOCTOR\\[-1pt]{\tiny see Module 6}};
\node[idrel] (HAS_BED) at (3.4,-7.8) {HAS\_BED};
\draw[partial] (HAS_BED) -- node[card,pos=0.72] {1} (WARD);
\draw[total] (HAS_BED) -- node[card,pos=0.72] {N} (BED);
\node[rel] (IN_DEPT) at (7.3,-4.4) {IN\_DEPT};
\draw[partial] (IN_DEPT) -- node[card,pos=0.72] {1} (DEPARTMENT);
\draw[total] (IN_DEPT) -- node[card,pos=0.72] {N} (WARD);
\node[rel] (ADMITTED) at (9.6,-14.5) {ADMITTED};
\draw[partial] (ADMITTED) -- node[card,pos=0.72] {1} (PATIENT.east);
\draw[total] (ADMITTED) -- node[card,pos=0.72] {N} (ADMISSION.-120);
\node[rel] (ADMITS) at (14,-13.4) {ADMITS};
\draw[partial] (ADMITS) -- node[card,pos=0.72] {1} (DOCTOR);
\draw[total] (ADMITS) -- node[card,pos=0.72] {N} (ADMISSION.-30);
\node[rel] (ADMITTED_TO) at (11.2,-7.6) {ADMITTED\_TO};
\draw[partial] (ADMITTED_TO) -- node[card,pos=0.72] {1} (DEPARTMENT);
\draw[total] (ADMITTED_TO) -- node[card,pos=0.72] {N} (ADMISSION);
\node[rel] (ALLOCATED) at (7.3,-10) {ALLOCATED};
\draw[partial] (ALLOCATED) -- node[card,pos=0.72] {1} (BED.north east);
\draw[total] (ALLOCATED) -- node[card,pos=0.72] {N} (ADMISSION.north west);
\node[rel] (OCCUPIES) at (7.3,-12.6) {OCCUPIES};
\draw[partial] (OCCUPIES) -- node[card,pos=0.72] {N} (BED.south east);
\draw[total] (OCCUPIES) -- node[card,pos=0.72] {M} (ADMISSION.south west);
\draw[partial] (OCCUPIES) -- node[card,pos=0.72] {1} (PATIENT);
\node[att] (OCCUPIES-assigned_at) at (4.9,-13.8) {assigned\_at};
\draw (OCCUPIES-assigned_at) -- (OCCUPIES);
\node[att] (OCCUPIES-released_at) at (4.9,-14.9) {released\_at};
\draw (OCCUPIES-released_at) -- (OCCUPIES);
\node[att] (OCCUPIES-status) at (4.9,-12.7) {status};
\draw (OCCUPIES-status) -- (OCCUPIES);
\node[rel] (RESULTS_IN) at (12,-15) {RESULTS\_IN};
\draw[partial] (RESULTS_IN) -- node[card,pos=0.72] {1} (EMERGENCY_CASE.75);
\draw[partial] (RESULTS_IN) -- node[card,pos=0.72] {N} (ADMISSION.-75);
\node[rel] (ARRIVES) at (9,-17.2) {ARRIVES};
\draw[partial] (ARRIVES) -- node[card,pos=0.72] {1} (PATIENT);
\draw[partial] (ARRIVES) -- node[card,pos=0.72] {N} (EMERGENCY_CASE.north west);
\node[rel] (CARRIES) at (7.4,-21.2) {CARRIES};
\draw[partial] (CARRIES) -- node[card,pos=0.72] {1} (AMBULANCE);
\draw[partial] (CARRIES) -- node[card,pos=0.72] {N} (EMERGENCY_CASE.south west);
\node[rel] (ASSIGNED) at (14.2,-16.6) {ASSIGNED};
\draw[partial] (ASSIGNED) -- node[card,pos=0.72] {1} (DOCTOR);
\draw[partial] (ASSIGNED) -- node[card,pos=0.72] {N} (EMERGENCY_CASE.north east);
\node[rel] (GIVEN_BED) at (3.4,-16.4) {GIVEN\_BED};
\draw[partial] (GIVEN_BED) -- node[card,pos=0.72] {1} (BED.south);
\draw[partial] (GIVEN_BED) -- (6,-18.6) -- node[card,pos=0.72] {N} (EMERGENCY_CASE.west);
\node[rel] (COVERS) at (3.4,-2.4) {COVERS};
\draw[partial] (COVERS) -- node[card,pos=0.72] {1} (WARD);
\draw[total] (COVERS) -- node[card,pos=0.72] {N} (WARD_ASSIGNMENT);
\end{tikzpicture}
\minilegend
\caption{Module 5, Beds and admissions (page 1 of 2): wards, beds, admissions, emergencies and ambulances}\label{fig:m5a}
\end{figure}

\begin{figure}[p]\centering
% Generated by docs/report/gen/gen-er.js: diagram 5b (Module 5)
\begin{tikzpicture}[x=1cm,y=1cm]
\node[refent] (ADMISSION) at (8.5,-7.6) {ADMISSION\\[-1pt]{\tiny see Module 5, page 1}};
\node[ent] (BILL) at (8.5,-14.4) {BILL};
\node[att] (BILL-bill_id) at (11,-12.9) {\underline{bill\_id}};
\draw (BILL-bill_id) -- (BILL);
\node[att] (BILL-bill_date) at (12.2,-13.9) {bill\_date};
\draw (BILL-bill_date) -- (BILL);
\node[att,dashed] (BILL-total_amount) at (12.4,-15) {total\_amount};
\draw (BILL-total_amount) -- (BILL);
\node[att] (BILL-paid_amount) at (5.6,-15.6) {paid\_amount};
\draw (BILL-paid_amount) -- (BILL);
\node[att] (BILL-status) at (5.8,-13.2) {status};
\draw (BILL-status) -- (BILL);
\node[ent] (BILL_ITEM) at (8.5,-20.6) {BILL\_ITEM};
\node[att] (BILL_ITEM-item_id) at (5.6,-19.4) {\underline{item\_id}};
\draw (BILL_ITEM-item_id) -- (BILL_ITEM);
\node[att] (BILL_ITEM-description) at (5.4,-20.6) {description};
\draw (BILL_ITEM-description) -- (BILL_ITEM);
\node[att] (BILL_ITEM-category) at (5.8,-21.8) {category};
\draw (BILL_ITEM-category) -- (BILL_ITEM);
\node[att] (BILL_ITEM-quantity) at (11.4,-19.4) {quantity};
\draw (BILL_ITEM-quantity) -- (BILL_ITEM);
\node[att] (BILL_ITEM-unit_price) at (11.8,-20.6) {unit\_price};
\draw (BILL_ITEM-unit_price) -- (BILL_ITEM);
\node[att,dashed] (BILL_ITEM-total_price) at (11.4,-21.8) {total\_price};
\draw (BILL_ITEM-total_price) -- (BILL_ITEM);
\node[refent] (ENCOUNTER) at (1.6,-2.6) {ENCOUNTER\\[-1pt]{\tiny see Module 1}};
\node[refent] (VITAL_SIGN) at (1.6,-7.6) {VITAL\_SIGN\\[-1pt]{\tiny see Module 1}};
\node[refent] (CONSULTATION) at (1.6,-12.6) {CONSULTATION\\[-1pt]{\tiny see Module 1}};
\node[refent] (PRESCRIPTION) at (15.4,-2.6) {PRESCRIPTION\\[-1pt]{\tiny see Module 2}};
\node[refent] (MED_ADMINISTRATION) at (15.2,-7.6) {MED\_ADMINISTRATION\\[-1pt]{\tiny see Module 2}};
\node[refent] (LAB_ORDER) at (15.4,-12.6) {LAB\_ORDER\\[-1pt]{\tiny see Module 4}};
\node[refent] (PATIENT) at (15.4,-17.4) {PATIENT\\[-1pt]{\tiny see Module 1}};
\node[rel] (STAY_VISIT) at (4.8,-2.6) {STAY\_VISIT};
\draw[partial] (STAY_VISIT) -- node[card,pos=0.72] {1} (ADMISSION.north west);
\draw[partial] (STAY_VISIT) -- node[card,pos=0.72] {N} (ENCOUNTER);
\node[rel] (TAKEN_ON) at (4.8,-7.6) {TAKEN\_ON};
\draw[partial] (TAKEN_ON) -- node[card,pos=0.72] {1} (ADMISSION);
\draw[partial] (TAKEN_ON) -- node[card,pos=0.72] {N} (VITAL_SIGN);
\node[rel] (STAY_NOTE) at (4.8,-12.6) {STAY\_NOTE};
\draw[partial] (STAY_NOTE) -- node[card,pos=0.72] {1} (ADMISSION.south west);
\draw[partial] (STAY_NOTE) -- node[card,pos=0.72] {N} (CONSULTATION);
\node[rel] (STAY_RX) at (12.2,-2.6) {STAY\_RX};
\draw[partial] (STAY_RX) -- node[card,pos=0.72] {1} (ADMISSION.north east);
\draw[partial] (STAY_RX) -- node[card,pos=0.72] {N} (PRESCRIPTION);
\node[rel] (STAY_DOSE) at (12,-7.6) {STAY\_DOSE};
\draw[partial] (STAY_DOSE) -- node[card,pos=0.72] {1} (ADMISSION);
\draw[partial] (STAY_DOSE) -- node[card,pos=0.72] {N} (MED_ADMINISTRATION);
\node[rel] (STAY_LAB) at (12.2,-12.6) {STAY\_LAB};
\draw[partial] (STAY_LAB) -- node[card,pos=0.72] {1} (ADMISSION.south east);
\draw[partial] (STAY_LAB) -- node[card,pos=0.72] {N} (LAB_ORDER);
\node[rel] (BILLED) at (12.6,-17.4) {BILLED};
\draw[partial] (BILLED) -- node[card,pos=0.72] {1} (PATIENT);
\draw[total] (BILLED) -- (8.9,-17.4) -- node[card,pos=0.72] {N} (BILL.south);
\node[rel] (STAY_BILL) at (8.5,-11) {STAY\_BILL};
\draw[partial] (STAY_BILL) -- node[card,pos=0.72] {1} (ADMISSION);
\draw[partial] (STAY_BILL) -- node[card,pos=0.72] {N} (BILL);
\node[rel] (ITEMISES) at (8.5,-17.6) {ITEMISES};
\draw[partial] (ITEMISES) -- node[card,pos=0.72] {1} (BILL);
\draw[total] (ITEMISES) -- node[card,pos=0.72] {N} (BILL_ITEM);
\end{tikzpicture}
\minilegend
\caption{Module 5, Beds and admissions (page 2 of 2): what an admission covers, and billing}\label{fig:m5b}
\end{figure}

\begin{figure}[p]\centering
% Generated by docs/report/gen/gen-er.js: diagram 6a (Module 6)
\begin{tikzpicture}[x=1cm,y=1cm]
\node[ent] (USER) at (5,-3.4) {USER};
\node[att] (USER-user_id) at (2.4,-1.4) {\underline{user\_id}};
\draw (USER-user_id) -- (USER);
\node[att] (USER-username) at (4.2,-0.8) {username};
\draw (USER-username) -- (USER);
\node[att] (USER-role) at (6,-0.8) {role};
\draw (USER-role) -- (USER);
\node[att] (USER-full_name) at (7.6,-1.4) {full\_name};
\draw (USER-full_name) -- (USER);
\node[att] (USER-email) at (2.2,-4.6) {email};
\draw (USER-email) -- (USER);
\node[ent] (WARD_ASSIGNMENT) at (13.2,-3.4) {WARD\_ASSIGNMENT};
\node[att] (WARD_ASSIGNMENT-assignment_id) at (15.4,-1.4) {\underline{assignment\_id}};
\draw (WARD_ASSIGNMENT-assignment_id) -- (WARD_ASSIGNMENT);
\node[att] (WARD_ASSIGNMENT-start_date) at (12.4,-1) {start\_date};
\draw (WARD_ASSIGNMENT-start_date) -- (WARD_ASSIGNMENT);
\node[att] (WARD_ASSIGNMENT-end_date) at (15.6,-4.6) {end\_date};
\draw (WARD_ASSIGNMENT-end_date) -- (WARD_ASSIGNMENT);
\node[ent] (AUDIT_LOG) at (2.4,-10.6) {AUDIT\_LOG};
\node[att] (AUDIT_LOG-audit_id) at (0.9,-12.4) {\underline{audit\_id}};
\draw (AUDIT_LOG-audit_id) -- (AUDIT_LOG);
\node[att] (AUDIT_LOG-action) at (2.4,-13) {action};
\draw (AUDIT_LOG-action) -- (AUDIT_LOG);
\node[att] (AUDIT_LOG-outcome) at (4,-12.6) {outcome};
\draw (AUDIT_LOG-outcome) -- (AUDIT_LOG);
\node[att] (AUDIT_LOG-created_at) at (0.9,-8.8) {created\_at};
\draw (AUDIT_LOG-created_at) -- (AUDIT_LOG);
\node[att] (AUDIT_LOG-details) at (4.4,-9) {details};
\draw (AUDIT_LOG-details) -- (AUDIT_LOG);
\node[ent] (DOCTOR) at (7.8,-10.6) {DOCTOR};
\node[att] (DOCTOR-doctor_id) at (6,-8.6) {\underline{doctor\_id}};
\draw (DOCTOR-doctor_id) -- (DOCTOR);
\node[att] (DOCTOR-specialization) at (9.6,-8.6) {specialization};
\draw (DOCTOR-specialization) -- (DOCTOR);
\node[att] (DOCTOR-shift) at (5.2,-11.9) {shift};
\draw (DOCTOR-shift) -- (DOCTOR);
\node[att] (DOCTOR-max_workload) at (6.4,-12.9) {max\_workload};
\draw (DOCTOR-max_workload) -- (DOCTOR);
\node[att,dashed] (DOCTOR-current_workload) at (11.2,-10.6) {current\_workload};
\draw (DOCTOR-current_workload) -- (DOCTOR);
\node[ent] (DEPARTMENT) at (13.2,-14.6) {DEPARTMENT};
\node[att] (DEPARTMENT-department_id) at (15.6,-13.2) {\underline{department\_id}};
\draw (DEPARTMENT-department_id) -- (DEPARTMENT);
\node[att] (DEPARTMENT-name) at (15.8,-14.4) {name};
\draw (DEPARTMENT-name) -- (DEPARTMENT);
\node[att] (DEPARTMENT-floor) at (15.6,-15.6) {floor};
\draw (DEPARTMENT-floor) -- (DEPARTMENT);
\node[att] (DEPARTMENT-phone) at (13.6,-16.6) {phone};
\draw (DEPARTMENT-phone) -- (DEPARTMENT);
\node[weakent] (DOCTOR_SCHEDULE) at (3.6,-18.4) {DOCTOR\_SCHEDULE};
\node[att] (DOCTOR_SCHEDULE-day_of_week) at (1.4,-17) {\dashuline{day\_of\_week}};
\draw (DOCTOR_SCHEDULE-day_of_week) -- (DOCTOR_SCHEDULE);
\node[att] (DOCTOR_SCHEDULE-start_time) at (1.2,-18.4) {\dashuline{start\_time}};
\draw (DOCTOR_SCHEDULE-start_time) -- (DOCTOR_SCHEDULE);
\node[att] (DOCTOR_SCHEDULE-end_time) at (1.4,-19.8) {end\_time};
\draw (DOCTOR_SCHEDULE-end_time) -- (DOCTOR_SCHEDULE);
\node[att] (DOCTOR_SCHEDULE-is_available) at (3.6,-20.4) {is\_available};
\draw (DOCTOR_SCHEDULE-is_available) -- (DOCTOR_SCHEDULE);
\node[weakent] (DAILY_ANALYTICS) at (8.6,-20.6) {DAILY\_ANALYTICS};
\node[att] (DAILY_ANALYTICS-analytics_date) at (6.4,-22) {\dashuline{analytics\_date}};
\draw (DAILY_ANALYTICS-analytics_date) -- (DAILY_ANALYTICS);
\node[att] (DAILY_ANALYTICS-patients_seen) at (8.6,-22.6) {patients\_seen};
\draw (DAILY_ANALYTICS-patients_seen) -- (DAILY_ANALYTICS);
\node[att] (DAILY_ANALYTICS-cases_pending) at (10.8,-22) {cases\_pending};
\draw (DAILY_ANALYTICS-cases_pending) -- (DAILY_ANALYTICS);
\node[refent] (WARD) at (13.2,-9) {WARD\\[-1pt]{\tiny see Module 5}};
\node[refent] (LAB_TEST) at (15.4,-20.4) {LAB\_TEST\\[-1pt]{\tiny see Module 4}};
\node[refent] (ADMISSION) at (11.6,-20.4) {ADMISSION\\[-1pt]{\tiny see Module 5}};
\node[rel] (OFFERS) at (15.4,-17.6) {OFFERS};
\draw[partial] (OFFERS) -- node[card,pos=0.72] {1} (DEPARTMENT.south east);
\draw[partial] (OFFERS) -- node[card,pos=0.72] {N} (LAB_TEST);
\node[rel] (IN_DEPT) at (13.2,-11.8) {IN\_DEPT};
\draw[partial] (IN_DEPT) -- node[card,pos=0.72] {1} (DEPARTMENT);
\draw[total] (IN_DEPT) -- node[card,pos=0.72] {N} (WARD);
\node[rel] (ADMITTED_TO) at (11.6,-17.6) {ADMITTED\_TO};
\draw[partial] (ADMITTED_TO) -- node[card,pos=0.72] {1} (DEPARTMENT.south west);
\draw[total] (ADMITTED_TO) -- node[card,pos=0.72] {N} (ADMISSION);
\node[rel] (IS_DOCTOR) at (6.6,-6.8) {IS\_DOCTOR};
\draw[partial] (IS_DOCTOR) -- node[card,pos=0.72] {1} (USER.south);
\draw[total] (IS_DOCTOR) -- node[card,pos=0.72] {1} (DOCTOR);
\node[rel] (WORKS_IN) at (10.6,-13.6) {WORKS\_IN};
\draw[partial] (WORKS_IN) -- node[card,pos=0.72] {1} (DEPARTMENT.west);
\draw[total] (WORKS_IN) -- node[card,pos=0.72] {N} (DOCTOR.south east);
\node[idrel] (HAS_SLOT) at (5.4,-15.4) {HAS\_SLOT};
\draw[partial] (HAS_SLOT) -- node[card,pos=0.72] {1} (DOCTOR.south west);
\draw[total] (HAS_SLOT) -- node[card,pos=0.72] {N} (DOCTOR_SCHEDULE);
\node[idrel] (HAS_STATS) at (8.4,-16.6) {HAS\_STATS};
\draw[partial] (HAS_STATS) -- node[card,pos=0.72] {1} (DOCTOR.south);
\draw[total] (HAS_STATS) -- node[card,pos=0.72] {N} (DAILY_ANALYTICS);
\node[rel] (NURSE_OF) at (9.1,-2.4) {NURSE\_OF};
\draw[partial] (NURSE_OF) -- node[card,pos=0.72] {1} (USER.north east);
\draw[total] (NURSE_OF) -- node[card,pos=0.72] {N} (WARD_ASSIGNMENT.north west);
\node[rel] (ASSIGNS) at (9.1,-4.4) {ASSIGNS};
\draw[partial] (ASSIGNS) -- node[card,pos=0.72] {1} (USER.south east);
\draw[partial] (ASSIGNS) -- node[card,pos=0.72] {N} (WARD_ASSIGNMENT.south west);
\node[rel] (COVERS) at (13.2,-6.2) {COVERS};
\draw[partial] (COVERS) -- node[card,pos=0.72] {1} (WARD);
\draw[total] (COVERS) -- node[card,pos=0.72] {N} (WARD_ASSIGNMENT);
\node[rel] (LOGS) at (2.4,-6.8) {LOGS};
\draw[partial] (LOGS) -- node[card,pos=0.72] {1} (USER.south west);
\draw[partial] (LOGS) -- node[card,pos=0.72] {N} (AUDIT_LOG);
\end{tikzpicture}
\minilegend
\caption{Module 6, Staff and audit (page 1 of 2): accounts, doctors, departments, assignments, schedules and the audit log}\label{fig:m6a}
\end{figure}

\begin{figure}[p]\centering
% Generated by docs/report/gen/gen-er.js: diagram 6b (Module 6)
\begin{tikzpicture}[x=1cm,y=1cm]
\node[refent] (DOCTOR) at (6.2,-11.6) {DOCTOR\\[-1pt]{\tiny see Module 6, page 1}};
\node[refent] (USER) at (10.8,-11.6) {USER\\[-1pt]{\tiny see Module 6, page 1}};
\node[refent] (ENCOUNTER) at (8.5,-1.6) {ENCOUNTER\\[-1pt]{\tiny see Module 1}};
\node[refent] (APPOINTMENT) at (3.28,-5.77) {APPOINTMENT\\[-1pt]{\tiny see Module 1}};
\node[refent] (CONSULTATION) at (2.29,-7.39) {CONSULTATION\\[-1pt]{\tiny see Module 1}};
\node[refent] (PRESCRIPTION) at (1.67,-9.39) {PRESCRIPTION\\[-1pt]{\tiny see Module 2}};
\node[refent] (SURGERY_REQUEST) at (1.45,-11.59) {SURGERY\_REQUEST\\[-1pt]{\tiny see Module 2}};
\node[refent] (LAB_ORDER) at (1.66,-13.79) {LAB\_ORDER\\[-1pt]{\tiny see Module 4}};
\node[refent] (ADMISSION) at (2.29,-15.79) {ADMISSION\\[-1pt]{\tiny see Module 5}};
\node[refent] (EMERGENCY_CASE) at (3.26,-17.42) {EMERGENCY\_CASE\\[-1pt]{\tiny see Module 5}};
\node[refent] (VITAL_SIGN) at (13.76,-6.08) {VITAL\_SIGN\\[-1pt]{\tiny see Module 1}};
\node[refent] (MED_ADMINISTRATION) at (14.78,-8) {MED\_ADMINISTRATION\\[-1pt]{\tiny see Module 2}};
\node[refent] (MEDICINE_BATCH) at (15.33,-10.35) {MEDICINE\_BATCH\\[-1pt]{\tiny see Module 3}};
\node[refent] (STOCK_MOVEMENT) at (15.33,-12.85) {STOCK\_MOVEMENT\\[-1pt]{\tiny see Module 3}};
\node[refent] (LAB_RESULT) at (14.78,-15.2) {LAB\_RESULT\\[-1pt]{\tiny see Module 4}};
\node[refent] (LAB_ATTACHMENT) at (13.76,-17.12) {LAB\_ATTACHMENT\\[-1pt]{\tiny see Module 4}};
\node[rel] (ATTENDS) at (4.66,-8.53) {ATTENDS};
\draw[partial] (ATTENDS) -- node[card,pos=0.72] {1} (DOCTOR);
\draw[total] (ATTENDS) -- node[card,pos=0.72] {N} (APPOINTMENT);
\node[rel] (CONDUCTS) at (6.6,-5) {CONDUCTS};
\draw[partial] (CONDUCTS) -- node[card,pos=0.72] {1} (DOCTOR);
\draw[total] (CONDUCTS) -- node[card,pos=0.72] {N} (ENCOUNTER);
\node[rel] (OPENS) at (10.4,-5) {OPENS};
\draw[partial] (OPENS) -- node[card,pos=0.72] {1} (USER);
\draw[partial] (OPENS) -- node[card,pos=0.72] {N} (ENCOUNTER);
\node[rel] (RECORDS) at (12.41,-8.61) {RECORDS};
\draw[partial] (RECORDS) -- node[card,pos=0.72] {1} (USER);
\draw[total] (RECORDS) -- node[card,pos=0.72] {N} (VITAL_SIGN);
\node[rel] (WRITES) at (4.14,-9.38) {WRITES};
\draw[partial] (WRITES) -- node[card,pos=0.72] {1} (DOCTOR);
\draw[total] (WRITES) -- node[card,pos=0.72] {N} (CONSULTATION);
\node[rel] (PRESCRIBES) at (3.81,-10.43) {PRESCRIBES};
\draw[partial] (PRESCRIBES) -- node[card,pos=0.72] {1} (DOCTOR);
\draw[total] (PRESCRIBES) -- node[card,pos=0.72] {N} (PRESCRIPTION);
\node[rel] (GIVEN_BY) at (12.97,-9.65) {GIVEN\_BY};
\draw[partial] (GIVEN_BY) -- node[card,pos=0.72] {1} (USER);
\draw[partial] (GIVEN_BY) -- node[card,pos=0.72] {N} (MED_ADMINISTRATION);
\node[rel] (REQUESTS) at (3.7,-11.59) {REQUESTS};
\draw[partial] (REQUESTS) -- node[card,pos=0.72] {1} (DOCTOR);
\draw[total] (REQUESTS) -- node[card,pos=0.72] {N} (SURGERY_REQUEST);
\node[rel] (RECEIVES) at (13.26,-10.92) {RECEIVES};
\draw[partial] (RECEIVES) -- node[card,pos=0.72] {1} (USER);
\draw[partial] (RECEIVES) -- node[card,pos=0.72] {N} (MEDICINE_BATCH);
\node[rel] (PERFORMS) at (13.26,-12.28) {PERFORMS};
\draw[partial] (PERFORMS) -- node[card,pos=0.72] {1} (USER);
\draw[total] (PERFORMS) -- node[card,pos=0.72] {N} (STOCK_MOVEMENT);
\node[rel] (ORDERS) at (3.81,-12.75) {ORDERS};
\draw[partial] (ORDERS) -- node[card,pos=0.72] {1} (DOCTOR);
\draw[total] (ORDERS) -- node[card,pos=0.72] {N} (LAB_ORDER);
\node[rel] (ANALYSES) at (12.97,-13.55) {ANALYSES};
\draw[partial] (ANALYSES) -- node[card,pos=0.72] {1} (USER);
\draw[partial] (ANALYSES) -- node[card,pos=0.72] {N} (LAB_RESULT);
\node[rel] (UPLOADS) at (12.41,-14.59) {UPLOADS};
\draw[partial] (UPLOADS) -- node[card,pos=0.72] {1} (USER);
\draw[total] (UPLOADS) -- node[card,pos=0.72] {N} (LAB_ATTACHMENT);
\node[rel] (ADMITS) at (4.14,-13.81) {ADMITS};
\draw[partial] (ADMITS) -- node[card,pos=0.72] {1} (DOCTOR);
\draw[total] (ADMITS) -- node[card,pos=0.72] {N} (ADMISSION);
\node[rel] (ASSIGNED) at (4.65,-14.66) {ASSIGNED};
\draw[partial] (ASSIGNED) -- node[card,pos=0.72] {1} (DOCTOR);
\draw[partial] (ASSIGNED) -- node[card,pos=0.72] {N} (EMERGENCY_CASE);
\end{tikzpicture}
\minilegend
\caption{Module 6, Staff and audit (page 2 of 2): the records each doctor and each staff account is linked to}\label{fig:m6b}
\end{figure}

\section{Explanation of the ER Diagram}
\textbf{Reading the diagrams.} A line ends at an entity with the cardinality that entity has in the
relationship. For example, \textsc{books} reads ``one PATIENT books N APPOINTMENTs''. A double line means
every instance of that entity takes part: every appointment is for a patient (\tbl{appointments.patient\_id}
is NOT NULL), while a patient need not have an appointment (single line). Optional links, such as a
prescription written outside a visit, are partial on both sides.

\textbf{Weak entities.} An entity is weak here only when the schema identifies it within its owner, that is,
its natural identifier is unique only together with the owner's key and the owner is mandatory:
\begin{itemize}
\item \textbf{BED} (owner WARD, partial key \dashuline{bed\_number}): bed numbers are unique only within a ward, UNIQUE (\tbl{ward\_id, bed\_number}).
\item \textbf{MEDICINE\_BATCH} (owner MEDICINE, partial key \dashuline{batch\_number}): a supplier's batch number is unique only for one medicine, UNIQUE (\tbl{medicine\_id, batch\_number}).
\item \textbf{DOCTOR\_SCHEDULE} (owner DOCTOR, partial key \dashuline{day\_of\_week}, \dashuline{start\_time}): a slot exists only as a slot of one doctor, UNIQUE (\tbl{doctor\_id, day\_of\_week, start\_time}), deleted with the doctor (ON DELETE CASCADE).
\item \textbf{DAILY\_ANALYTICS} (owner DOCTOR, partial key \dashuline{analytics\_date}): one statistics row per doctor per day, UNIQUE (\tbl{doctor\_id, analytics\_date}), deleted with the doctor.
\end{itemize}
In the tables each also has a surrogate integer key, added during mapping (Section~4.2). Line items such as
BILL\_ITEM are \emph{not} modelled as weak: nothing in the schema identifies them within their bill, so their own
key identifies them.

\textbf{Module 1, Patients and visits} (Figures~\ref{fig:m1a}, \ref{fig:m1b}). PATIENT books APPOINTMENTs and has
ENCOUNTERs (visits). An appointment \textsc{becomes} at most one encounter and an encounter comes from at most one
appointment (1:1, both partial: walk-ins and emergencies have no appointment). VITAL\_SIGN readings and
CONSULTATION notes always name their patient (total) and are filed under a visit (\textsc{taken\_in},
\textsc{note\_in}) or, for inpatients, an admission (\textsc{taken\_on}, \textsc{stay\_note}). A consultation may
also refer to its appointment (\textsc{note\_for}). Attributes show the derived \emph{age} (computed from
\emph{date\_of\_birth}, never stored), the multivalued \emph{allergies}, and the composite \emph{slot},
\emph{timeline} and \emph{blood\_pressure}. Links to other modules: DOCTOR \textsc{attends} appointments,
\textsc{conducts} encounters and \textsc{writes} notes; USER \textsc{opens} an encounter and \textsc{records} vitals
(Module~6); PATIENT is \textsc{admitted}, \textsc{arrives} as an emergency, is \textsc{billed}, and appears in
\textsc{occupies} (Module~5); PATIENT, ENCOUNTER and CONSULTATION are linked to prescriptions and surgery requests
(Module~2) and lab orders (Module~4).

\textbf{Module 2, Clinical orders} (Figure~\ref{fig:m2}). A PRESCRIPTION \textsc{lists} its PRESCRIPTION\_ITEMs
(total on the item side). Each item is \textsc{of\_drug} one MEDICINE and is \textsc{filled\_from} one or more
MEDICINE\_BATCHes (M:N, with the relationship attribute \emph{quantity}). Every MED\_ADMINISTRATION (a dose on the
MAR) is \textsc{given\_as} part of one item (total) and optionally belongs to an admission and to the nurse who
recorded it. A prescription and a SURGERY\_REQUEST each belong to a patient and a doctor (total) and optionally to
the visit, note or admission they were written in.

\textbf{Module 3, Pharmacy} (Figure~\ref{fig:m3}). MEDICINE \textsc{has\_batch} MEDICINE\_BATCHes through an
identifying relationship. Every STOCK\_MOVEMENT \textsc{moves} one medicine (total) and, when it touched a batch,
comes \textsc{from\_batch} it; movements from before batches existed have none (partial). The medicine's
\emph{stock\_quantity} is drawn as derived: it is the sum of its batches, stored for speed and kept equal in the
same transaction (Section~6.6). USER links: \textsc{receives} (batches) and \textsc{performs} (movements).

\textbf{Module 4, Laboratory} (Figure~\ref{fig:m4}). A LAB\_ORDER is \textsc{of\_test} one LAB\_TEST, and
\textsc{yields} at most one LAB\_RESULT (1:1; every result has its order, total). A result has any number of
LAB\_ATTACHMENTs. Orders link to the patient and doctor (total) and optionally to the visit, note and admission;
tests to the DEPARTMENT that \textsc{offers} them; results and files to the USER who \textsc{analyses} and
\textsc{uploads} them.

\textbf{Module 5, Beds and admissions} (Figures~\ref{fig:m5a}, \ref{fig:m5b}). A WARD \textsc{has\_bed} BEDs
(identifying). An ADMISSION is \textsc{allocated} one bed, is \textsc{admitted} for one patient by one doctor
(\textsc{admits}) to one department (\textsc{admitted\_to}), all total. \textsc{occupies} is ternary: each row of
the bed history joins one bed, one admission and its patient, with the attributes \emph{assigned\_at},
\emph{released\_at} and \emph{status}; every admission takes part (total), because a trigger writes the row.
An EMERGENCY\_CASE may be \textsc{carried} by an AMBULANCE, \textsc{assigned} a doctor and \textsc{given\_bed}, and
\textsc{results\_in} an admission; all partial, because cases wait in a queue first. Figure~\ref{fig:m5b} shows the
six records an admission can cover (Modules 1, 2 and 4) and billing: a BILL is \textsc{billed} to one patient
(total), optionally for an admission, and \textsc{itemises} BILL\_ITEMs whose \emph{total\_price} is derived
(quantity $\times$ unit price); the bill's \emph{total\_amount} is the stored sum of its items.

\textbf{Module 6, Staff and audit} (Figures~\ref{fig:m6a}, \ref{fig:m6b}). A USER account \textsc{is\_doctor} for at
most one DOCTOR (1:1; every doctor has an account). A doctor \textsc{works\_in} one DEPARTMENT and has
DOCTOR\_SCHEDULE slots and DAILY\_ANALYTICS rows (both weak). A WARD\_ASSIGNMENT joins a nurse
(\textsc{nurse\_of}, total) to a ward (\textsc{covers}, total) for a period of days, ended early at the moment an
administrator ends it (\emph{ended\_at}), and records the administrator who \textsc{assigns} it. Every AUDIT\_LOG row \textsc{logs} the user who acted, when known (failed logins have none).
Figure~\ref{fig:m6b} gathers, for DOCTOR and USER, every relationship that enters them from Modules~1 to~5, so the
two pages of Module~6 together show all 15 links of these two entities.
